% from ega2-1-fr.tex
\oldpage[II]{7}
de $S$ (voir (\hyperref[II.1.3.3-fr]{1.3.3})). Toutefois, rappelons que si
$S$ est un \emph{schéma}, tout $S$-préschéma qui est un schéma affine est
affine au-dessus de $S$ (\hyperref[I.5.5.10-fr]{I, 5.5.10}).
\end{remark}

\begin{proposition}[1.2.4]
\label{II.1.2.4-fr}
Tout $S$-préschéma qui est affine au-dessus de $S$ est séparé au-dessus de
$S$ (autrement dit, est un $S$-schéma).
\end{proposition}

Cela résulte aussitôt de (\hyperref[I.5.5.5-fr]{I, 5.5.5}) et
(\hyperref[I.5.5.8-fr]{I, 5.5.8}).

\begin{proposition}[1.2.5]
\label{II.1.2.5-fr}
Soient $X$ un $S$-schéma affine au-dessus de $S$, $f:X\to S$ le morphisme
structural. Pour tout ouvert $U\subset S$, $f^{-1}(U)$ est affine au-dessus
de $U$.
\end{proposition}

En vertu de la déf. (\hyperref[II.1.2.1-fr]{1.2.1}), on est ramené au cas où
$S=\operatorname{Spec}(A)$ et $X=\operatorname{Spec}(B)$ sont affines et
alors $f=({}^a\varphi,\widetilde{\varphi})$, où $\varphi:A\to B$ est un
homomorphisme. Comme les $D(g)$, où $g\in A$, forment une base de $S$, on est
ramené au cas où $U=D(g)$~; mais on sait alors que
$f^{-1}(U)=D(\varphi(g))$ (\hyperref[I.1.2.2.2-fr]{I, 1.2.2.2}), d'où la
proposition.

\begin{proposition}[1.2.6]
\label{II.1.2.6-fr}
Soient $X$ un $S$-schéma affine au-dessus de $S$, $f:X\to S$ le morphisme
structural. Pour tout $\mathcal{O}_X$-Module quasi-cohérent $\mathcal{F}$,
$f_*(\mathcal{F})$ est un $\mathcal{O}_S$-Module quasi-cohérent.
\end{proposition}

Compte tenu de (\hyperref[II.1.2.4-fr]{1.2.4}), cela résulte de
(\hyperref[I.9.2.2-fr]{I, 9.2.2, a)}).

En particulier, la $\mathcal{O}_S$-Algèbre
$\mathcal{A}(X)=f_*(\mathcal{O}_X)$ est \emph{quasi-cohérente}.

\begin{proposition}[1.2.7]
\label{II.1.2.7-fr}
Soient $X$ un $S$-schéma affine au-dessus de $S$. Pour tout $S$-préschéma
$Y$, l'application $h\mapsto\mathcal{A}(h)$ de l'ensemble
$\operatorname{Hom}_S(Y,X)$ dans l'ensemble
$\operatorname{Hom}(\mathcal{A}(X),\mathcal{A}(Y))$
(\hyperref[II.1.1.2-fr]{1.1.2}) est bijective.
\end{proposition}

Soient $f:X\to S$, $g:Y\to S$ les morphismes structuraux. Supposons d'abord
$S=\operatorname{Spec}(A)$ et $X=\operatorname{Spec}(B)$ affines~; il faut
prouver que pour tout homomorphisme
$\omega:f_*(\mathcal{O}_X)\to g_*(\mathcal{O}_Y)$ de
$\mathcal{O}_S$-Algèbres, il existe un $S$-morphisme $h:Y\to X$ et un seul
tel que $\mathcal{A}(h)=\omega$. Par définition, pour tout ouvert $U\subset
S$, $\omega$ définit un homomorphisme
\[
  \omega_U=\Gamma(U,\omega):
  \Gamma(f^{-1}(U),\mathcal{O}_X)\to
  \Gamma(g^{-1}(U),\mathcal{O}_Y)
\]
de $\Gamma(U,\mathcal{O}_S)$-algèbres. En particulier, pour $U=S$, cela
donne un homomorphisme
$\varphi:\Gamma(X,\mathcal{O}_X)\to\Gamma(Y,\mathcal{O}_Y)$ de
$\Gamma(S,\mathcal{O}_S)$-algèbres, auquel correspond un $S$-morphisme
$h:Y\to X$ bien déterminé, puisque $X$ est affine
(\hyperref[I.2.2.4-fr]{I, 2.2.4}). Reste à prouver que
$\mathcal{A}(h)=\omega$, autrement dit, que pour tout ouvert $U$ d'une base
de $S$, $\omega_U$ coïncide avec l'homomorphisme d'algèbres $\varphi_U$
correspondant au $S$-morphisme $g^{-1}(U)\to f^{-1}(U)$ restriction de $h$.
On peut se borner au cas où $U=D(\lambda)$, avec $\lambda\in S$~; alors, si
$f=({}^a\rho,\widetilde{\rho})$, où $\rho:A\to B$ est un homomorphisme
d'anneaux, on a $f^{-1}(U)=D(\mu)$, où $\mu=\rho(\lambda)$, et
$\Gamma(f^{-1}(U),\mathcal{O}_X)$ est l'anneau de fractions $B_\mu$~; or,
le diagramme
\[
  \xymatrix{
    B\ar[r]^\varphi\ar[d] & \Gamma(Y,\mathcal{O}_Y)\ar[d]\\
    B_\mu\ar[r]_{\varphi_U} & \Gamma(g^{-1}(U),\mathcal{O}_Y)
  }
\]
est commutatif, et il en est de même du diagramme analogue où $\varphi_U$
est remplacé par $\omega_U$~; l'égalité $\varphi_U=\omega_U$ résulte alors
de la propriété universelle des anneaux de fractions
(\hyperref[0.1.2.4-fr]{0, 1.2.4}).

Passons au cas général, et soit $(S_\alpha)$ un recouvrement de $S$ par des
ouverts affines
